% GENERATED FILE. Edit canonical lessons, then run npm run generate:books.
% canonical-source-hash: fnv1a64:ae19a944a36f33fc
% canonical-lessons: UR-C162-maqami, UR-C162-angrezi, UR-C162-pakistani, UR-C162-hindustani, UR-C162-surmai

\chapter{Local, English, Pakistani, Indian, Grey}
\label{ch:ur-local-english-pakistani-indian-grey}
\hlchaptermodality{162}

\begin{chapteropening}
\textbf{By the end of this chapter:} \emph{I can say local, English, Pakistani, Indian and grey.}

Complete the last lesson of chapter 162: I can say local, English, Pakistani, Indian and grey.
\end{chapteropening}

\section[\texorpdfstring{\ur{مقامی} (maqāmī)}{مقامی (maqāmī)}]{\ur{مقامی} (maqāmī) — local}

\label{lesson:UR-C162-maqami}



\begin{quote}
Before the new one: say the Urdu for rough, then the Urdu for round.
\end{quote}

\subsection*{You'll want to know: \ur{مقامی}}
\textbf{\ur{مقامی}} — \emph{maqāmī} — \textquotedblleft{}local\textquotedblright{}.

\begin{cousinweb}
\textbf{\ur{مقامی}} (\emph{maqāmī}) is \emph{maqām}, \textquotedblleft{}place\textquotedblright{}, with the ending \emph{-ī}, \textquotedblleft{}of\textquotedblright{}. \emph{maqām} is Arabic, \textquotedblleft{}a place of standing\textquotedblright{}, and reached Urdu through Persian (Lane, root \emph{q-w-m}).
\end{cousinweb}

\begin{grammarlens}[title={What you've built}]
One more word for this chapter.
\end{grammarlens}

\subsection*{Your turn}
\begin{itemize}
\raggedright
  \item \emph{Say it:} \emph{maqāmī}
  \item \emph{Say it:} \emph{maqāmī}, once more
  \item \emph{Say it:} say \emph{maqāmī}
  \item \emph{Recall:} say the Urdu for rough, then the Urdu for round, then say \emph{maqāmī} again
\end{itemize}

\begin{tcolorbox}[breakable,colback=teal!4,colframe=teal!35!black,title={Before you move on}]
What does \ur{مقامی} mean? (\textquotedblleft{}Local\textquotedblright{}.) Say it once more.
\end{tcolorbox}
\section[\texorpdfstring{\ur{انگریزی} (angrezī)}{انگریزی (angrezī)}]{\ur{انگریزی} (angrezī) — English}

\label{lesson:UR-C162-angrezi}



\begin{quote}
Before the new one: say the Urdu for round, then the Urdu for local.
\end{quote}

\subsection*{You'll want to know: \ur{انگریزی}}
\textbf{\ur{انگریزی}} — \emph{angrezī} — \textquotedblleft{}English\textquotedblright{}.

\begin{cousinweb}
\textbf{\ur{انگریزی}} (\emph{angrezī}) is \emph{angrez}, \textquotedblleft{}an Englishman\textquotedblright{}, with the ending \emph{-ī}. \emph{angrez} is usually traced, through Persian, to Portuguese \emph{inglês}, \textquotedblleft{}English\textquotedblright{}.
\end{cousinweb}

\begin{grammarlens}[title={What you've built}]
One more word for this chapter.
\end{grammarlens}

\subsection*{Your turn}
\begin{itemize}
\raggedright
  \item \emph{Say it:} \emph{angrezī}
  \item \emph{Say it:} \emph{angrezī}, once more
  \item \emph{Say it:} say \emph{angrezī}
  \item \emph{Recall:} say the Urdu for round, then the Urdu for local, then say \emph{angrezī} again
\end{itemize}

\begin{tcolorbox}[breakable,colback=teal!4,colframe=teal!35!black,title={Before you move on}]
What does \ur{انگریزی} mean? (\textquotedblleft{}English\textquotedblright{}.) Say it once more.
\end{tcolorbox}
\section[\texorpdfstring{\ur{پاکستانی} (pākistānī)}{پاکستانی (pākistānī)}]{\ur{پاکستانی} (pākistānī) — Pakistani}

\label{lesson:UR-C162-pakistani}



\begin{quote}
Before the new one: say the Urdu for local, then the Urdu for English.
\end{quote}

\subsection*{You'll want to know: \ur{پاکستانی}}
\textbf{\ur{پاکستانی}} — \emph{pākistānī} — \textquotedblleft{}Pakistani\textquotedblright{}.

\begin{cousinweb}
\textbf{\ur{پاکستانی}} (\emph{pākistānī}) is \emph{Pākistān} with the ending \emph{-ī}, \textquotedblleft{}of\textquotedblright{}. The name was coined in 1933, and is read as Persian \emph{pāk}, \textquotedblleft{}pure\textquotedblright{}, and \emph{-stān}, \textquotedblleft{}land\textquotedblright{}.
\end{cousinweb}

\begin{grammarlens}[title={What you've built}]
One more word for this chapter.
\end{grammarlens}

\subsection*{Your turn}
\begin{itemize}
\raggedright
  \item \emph{Say it:} \emph{pākistānī}
  \item \emph{Say it:} \emph{pākistānī}, once more
  \item \emph{Say it:} say \emph{pākistānī}
  \item \emph{Recall:} say the Urdu for local, then the Urdu for English, then say \emph{pākistānī} again
\end{itemize}

\begin{tcolorbox}[breakable,colback=teal!4,colframe=teal!35!black,title={Before you move on}]
What does \ur{پاکستانی} mean? (\textquotedblleft{}Pakistani\textquotedblright{}.) Say it once more.
\end{tcolorbox}
\section[\texorpdfstring{\ur{ہندوستانی} (hindustānī)}{ہندوستانی (hindustānī)}]{\ur{ہندوستانی} (hindustānī) — Indian}

\label{lesson:UR-C162-hindustani}



\begin{quote}
Before the new one: say the Urdu for English, then the Urdu for Pakistani.
\end{quote}

\subsection*{You'll want to know: \ur{ہندوستانی}}
\textbf{\ur{ہندوستانی}} — \emph{hindustānī} — \textquotedblleft{}Indian\textquotedblright{}.

\begin{cousinweb}
\textbf{\ur{ہندوستانی}} (\emph{hindustānī}) is \emph{Hindūstān}, \textquotedblleft{}the land of Hind\textquotedblright{}, with the ending \emph{-ī}, \textquotedblleft{}of\textquotedblright{}. The name is Persian: \emph{Hindū}, \textquotedblleft{}of Hind, India\textquotedblright{}, and \emph{-stān}, \textquotedblleft{}land\textquotedblright{}.
\end{cousinweb}

\begin{grammarlens}[title={What you've built}]
One more word for this chapter.
\end{grammarlens}

\subsection*{Your turn}
\begin{itemize}
\raggedright
  \item \emph{Say it:} \emph{hindustānī}
  \item \emph{Say it:} \emph{hindustānī}, once more
  \item \emph{Say it:} say \emph{hindustānī}
  \item \emph{Recall:} say the Urdu for English, then the Urdu for Pakistani, then say \emph{hindustānī} again
\end{itemize}

\begin{tcolorbox}[breakable,colback=teal!4,colframe=teal!35!black,title={Before you move on}]
What does \ur{ہندوستانی} mean? (\textquotedblleft{}Indian\textquotedblright{}.) Say it once more.
\end{tcolorbox}
\section[\texorpdfstring{\ur{سرمئی} (surmaī)}{سرمئی (surmaī)}]{\ur{سرمئی} (surmaī) — grey}

\label{lesson:UR-C162-surmai}



\begin{quote}
Before the new one: say the Urdu for Pakistani, then the Urdu for Indian.
\end{quote}

\subsection*{You'll want to know: \ur{سرمئی}}
\textbf{\ur{سرمئی}} — \emph{surmaī} — \textquotedblleft{}grey\textquotedblright{}.

\begin{cousinweb}
\textbf{\ur{سرمئی}} (\emph{surmaī}) is \textquotedblleft{}the colour of \emph{surma}\textquotedblright{}, the dark grey powder put round the eyes, with the ending \emph{-ī}. \emph{surma} is usually traced to Persian \emph{surma}.
\end{cousinweb}

\begin{grammarlens}[title={What you've built}]
One more word for this chapter.
\end{grammarlens}

\subsection*{Your turn}
\begin{itemize}
\raggedright
  \item \emph{Say it:} \emph{surmaī}
  \item \emph{Say it:} \emph{surmaī}, once more
  \item \emph{Say it:} say \emph{surmaī}
  \item \emph{Recall:} say the Urdu for Pakistani, then the Urdu for Indian, then say \emph{surmaī} again
\end{itemize}

\begin{tcolorbox}[breakable,colback=teal!4,colframe=teal!35!black,title={Before you move on}]
What does \ur{سرمئی} mean? (\textquotedblleft{}Grey\textquotedblright{}.) Say it once more.
\end{tcolorbox}
